\section{实战复盘二：安全与治理的工程化落地清单}

OpenClaw 引发的大部分争议，本质上是“高自治系统如何治理”。本章把访谈提到的安全理念转成可执行清单。

\subsection{权限分层：把能力切成风险等级}

把所有动作放在同一权限层，是最常见设计错误。应按风险分层并设置不同确认策略。

\begin{knowledgebox}{建议的四级权限模型}
\begin{itemize}
    \item \textbf{L1 读操作}：读取公开网页、知识库、非敏感文件；
    \item \textbf{L2 低风险写操作}：草稿、标签、提醒、临时文件；
    \item \textbf{L3 中风险执行}：发消息、提工单、改配置；
    \item \textbf{L4 高风险执行}：支付、删除、发布、外发敏感数据。
\end{itemize}
\end{knowledgebox}

其中 L3/L4 需要显式确认和操作证据留存。若系统支持“自动通过”，也应有时间窗口和撤回机制。

\subsection{模型策略：不是“最强模型”而是“最匹配模型”}

访谈提到“不要盲信便宜模型”并不意味着只用最贵模型，而是任务分层：
\begin{itemize}
    \item 低风险检索与草拟可用廉价模型；
    \item 高风险决策与外发动作使用高可靠模型；
    \item 关键操作前后做双模型一致性检查或规则复核。
\end{itemize}

\begin{warningbox}{低价模型的隐性成本}
若低价模型在高风险场景误判率上升，后续人工返工、安全事件、用户流失会吞噬全部成本优势。成本必须按“全链路”计算，而非按 token 单价计算。
\end{warningbox}

\subsection{可观测性：把“自治行为”变成“可追踪行为”}

\begin{importantbox}{必须记录的运行证据}
\begin{enumerate}
    \item 输入证据：消息、文件、网页来源；
    \item 推理证据：关键决策摘要与约束命中情况；
    \item 执行证据：工具调用参数、返回结果、错误码；
    \item 变更证据：配置差异、规则更新、回滚历史。
\end{enumerate}
\end{importantbox}

没有证据链，任何“agent 出错”都只能归因于“模型抽风”，这在生产环境不可接受。

\subsection{事件响应：把“模型事故”当成系统事故处理}

\begin{longtable}{p{0.2\textwidth}p{0.32\textwidth}p{0.38\textwidth}}
\toprule
\textbf{阶段} & \textbf{目标} & \textbf{动作} \\
\midrule
检测 & 尽快识别异常行为 & 告警规则、异常阈值、行为基线比较 \\
\midrule
隔离 & 阻断持续伤害 & 降级到只读模式、冻结高风险工具 \\
\midrule
定位 & 还原事故路径 & 回放日志、比对配置变更、重建上下文 \\
\midrule
修复 & 恢复服务稳定性 & 回滚策略、补丁发布、权限临时收紧 \\
\midrule
复盘 & 防止同类问题复发 & 增加测试样例、修订策略文档、更新 runbook \\
\bottomrule
\end{longtable}

\subsection{“安全焦虑”与“建设速度”如何平衡}

访谈给出的可行立场是：承认焦虑合理，但不陷入停摆。最务实路径是“持续公开风险、持续修复路径、持续缩小 blast radius”。

\begin{figure}[H]
\centering
\includegraphics[width=0.9\textwidth]{frames/frame-031210.jpg}
\caption{访谈收束段：社会风险讨论与用户真实收益并置\protect\footnotemark}
\end{figure}
\footnotetext{视频画面时间区间：03:12:05--03:12:20。}

\subsection{本章小结}
安全不是“发布前检查项”，而是 Agent 产品的日常运行能力。OpenClaw 的启发在于：越早把安全纳入产品主线，越能在扩散期保持控制力。

